\cotefile{76}
\scanpage{51}
\deleted{Ainsi, si $S$ est l'un des sommets du cube $C(B,\sigma)$ -- correspondant aux sections de $B$ sur $I=B/\sigma$, alors \unsure{tout} $s\in S$ est l'intersection des facettes}

\uline{Corollaire} Soient $(B,\sigma)$, $(B',\sigma')$ deux \deleted{\unsure{basi}} maquettes de cubes, \deleted{Alors} définissant des cubes $C$, $C'$, de sorte que l'on a
\[
\operatorname{Isom}((B,\sigma),(B',\sigma'))\xrightarrow{\alpha}
\operatorname{Isom}_{\mathrm{aff}}(C,C')\xrightarrow{\beta}
\operatorname{Isom}_{\mathrm{ord}}(\operatorname{Comb}(C),\operatorname{Comb}(C'))
\]
Ces applications sont des isomorphismes. En particulier on a
\[
\begin{aligned}
\operatorname{Aut}((B,\sigma))&\xrightarrow{\sim}\operatorname{Aut}_{\mathrm{aff}}(C)\xrightarrow{\sim}\operatorname{Aut}_{\mathrm{ord}}(\operatorname{Comb}(C))\\
\shortparallel\qquad&\\[-2pt]
\text{commutant de }\sigma\qquad&\\
\text{dans }\unsure{\mathfrak S_B|\mathscr S_B}\qquad&\\
\mathrel{\rotatebox{90}{$\simeq$}}\ \text{non can.}\qquad&\\
\mathfrak S_I\cdot\{\pm1\}^{I}\qquad&
\end{aligned}
\]

\uline{Dém} On sait déjà que $\beta\alpha$ est un iso (histoire combinatoire -- cf cor. précédents). Il reste : prouver que $\beta$ l'est, On \unsure{en} sait déjà que $\beta$ est surjectif ($\beta\alpha$ bij), or il est injectif (car un isom $C\unsure{-|\to} C'$ est connu quand on connaît son effet sur les seuls pts extrémaux) -- g.g.

\uline{NB} Au lieu de $\operatorname{Isom}_{\mathrm{aff}}(C,C')$ on \deleted{peut} aussi d'écrire $\operatorname{Isom}_{\mathrm{proj}}(C,C')$

% slanted
\margin{\unsure{remonté|remonter}}
\uline{Corollaire} Les facettes du cube sont des cubes
% archive sheet number: 50
